% Copyright (c) 2026 Eugene Joseph M. Ragasa.
% SPDX-License-Identifier: Apache-2.0
% Author: Eugene Joseph M. Ragasa
% Formal proof completed: 2026-10-28
\documentclass[11pt]{article}
\usepackage[T1]{fontenc}\usepackage{amsmath,amssymb}\usepackage[margin=1in]{geometry}\usepackage[hidelinks]{hyperref}
\title{Proof Dossier: Squared Matrix-Frobenius Fourier Truncation-Channel RMS}\author{Eugene Joseph M. Ragasa}\date{Formal proof completed: 2026-10-28}
\begin{document}\raggedright\maketitle
\section{Identity and status}
ID \texttt{BG-FOURIER-THM-073}; status \texttt{proved}. Lean declaration:
\par\noindent{\tiny\path{ProjectKoios.BandGauge.Fourier.finiteMatrixFrobeniusRootMeanSquare_truncationResidual_sq_eq_omitted}}.
\section{Statement and proof}
For complex matrix coefficients on a nonempty finite abelian group, the squared finite root-mean-square entrywise Frobenius norm of the deliberate-truncation channel equals the omitted coefficient Frobenius power. The proof squares \texttt{BG-FOURIER-THM-072} and proves nonnegativity entrywise before summing over modes, rows, and columns.
\section{Boundary and audit}
The equality is exact finite Frobenius algebra. It does not check serialized or binary64 matrix values, physical hopping range, or numerical adequacy. Axioms are \texttt{propext}, \texttt{Classical.choice}, and \texttt{Quot.sound}; no admissions. Lean \texttt{4.35.0}, Mathlib {\scriptsize\path{56a005855ea17e0b0958f73ec801ae140b0a1679}}; AI-assisted, kernel-checked, no independent expert review recorded.
\end{document}
